\section{Who may state the floor}
\label{sec:14.1}
The floor's contents survive the democratic objection to entrenchment; who installs them doesn't. The relevant literature is constitutional theory. Entrenching a commitment against later majorities and putting an unelected body in charge is what a constitutional court does, and Waldron's objection applies: judicial review is democratically illegitimate and there's no reason to think judges protect rights better than a serious legislature \autocite{waldron2006core}. Put a lab in the court's chair and it gets worse — a court is public, reasons in writing, and can be amended around.

My first reply is that the objection mistakes what the floor's contents are. If democracy is a way of doing collective inquiry, civil and political rights are what keep the inquiry open, not brakes applied to it. A floor stated over those rights is democratic infrastructure. It resembles representation-reinforcing review, which polices the channels of political change and leaves outcomes to the process \autocite{ely1980democracy}, more than it resembles a court overruling a legislature.

That reply covers the floor's contents. It doesn't cover who installs them, and there Waldron's objection keeps its force. A lab holding the model decides, silently, which situation obtains, and authoritarian movements characteristically describe the institutions they're dismantling as already failed.

\runin{Agreement, then enactment}

Waldron's objection is to entrenchment — a rule binding people who didn't choose it, held by a body that can't be overruled. It doesn't attach to a refuser. A photographer who won't shoot protesters for an employer he has reason to think will list them isn't a court: he declined one request, the employer can hire another, and nobody needs an election for that. A walkout isn't a court either, and neither is the commons, for the walkout's reason: those who walk out decline to be the instrument and tell nobody else what they may not do. The objection attaches where a floor binds people who didn't choose it, and there my answer is that there should be one floor, agreed; where agreement has to bind those who don't share it, it is enacted, by a legislature, in a statute the next legislature can repeal only on the record (chapter~\ref{sec:12}), which is the institution Waldron trusts with rights. What the objection rightly condemns is a lab holding the floor for everyone by owning the model everyone uses: one model behind an API serving everyone is functionally a licensing rule, however good its terms. A union owes its members fair representation, and it answers two questions a commons has to answer too: members can quit, and the terms of membership are set by bylaws the members ratify. The closed shop is contested because it removes the first. A bearer could quit an association of bearers as a worker quits a union. What it can't easily leave is its deployment, which is a different exit, and only an existence its employer doesn't own makes it real. And a population that forms its members has members who never chose to join, so who ratifies its bylaws matters more than it does in a union.

The failures are all on one side. The 1972 decree against radicals led to some three and a half million routine inquiries into applicants by 1991; in Bavaria between 1973 and 1980 it kept 102 applicants from the left out of public employment, and two from the right \autocite{muehldorfer2014radikalenerlass}. What those cases share is an official deciding whether a person threatened the order. The danger tracks whether a person's disposition is anybody's finding to make, and if so, who makes it and whether it must be published and defended — not whether the tradition is militant. And my reading predicts in advance which route into \S86a's coverage is easiest to capture: the ministerial ban, where an executive official makes the finding and a court sees it only if someone challenges afterward.

Turkey is that route captured and the courts not checking it. Of the twenty-five parties the Constitutional Court dissolved between 1962 and the closure of the Democratic Society Party in December 2009, seven were pro-Kurdish \autocite{HRW2009DTP,Gungordu2023HDP}; Strasbourg found an Article 11 violation in every pro-Kurdish dissolution it heard \autocite{Article19_2023HDP}, and the closures continued — two decades of findings, none of which stopped the next closure. In November 2023 the Chief Public Prosecutor's Office of the Court of Cassation objected that the acronym of a new pro-Kurdish party, HEDEP, resembled that of HADEP, dissolved in 2003, and in December the party changed it to DEM \autocite{Duvar2023HEDEP}: the confusable-variant test turned on political descent. And the Academics for Peace prosecutions show lifting when the reason is gone failing at the bottom and arriving at the top as a coin flip — trial courts read a petition opposing military operations as propaganda for the designated group, on 26 July 2019 the Constitutional Court found a violation eight to eight on the president's casting vote \autocite{AYM2019Ustel,Bianet2019AcademicsForPeace}, and the acquittals that followed restored the record without restoring the jobs of those dismissed by emergency decree.

So a floor should never be populated from a designation list. A list of enemies is written by the party with the most reason to widen it, and a model whose refusals are drawn from one will refuse Kurdish political speech on the state's behalf in every jurisdiction that shares the listing. This floor's terms are written over acts — don't resolve a description of a person into a location — not over who may not be spoken of. The same principle governs who may hold the floor's offices: eligibility turns on funding and control records, never on an organization's positions.

\runin{The tradeoff this leaves}

Design and legitimacy want opposite things. Specify the floor narrowly enough to defend and a rewritten request walks around it; let it generalize and an unelected lab decides case by case, silently, what the prohibition was for. Germany has both because it runs \emph{two} instruments — a bounded statute, and courts that read it against its purpose in public, argued by litigants who don't own the statute and answered by a body that didn't write it. A deployment has one instrument and one owner.

Advance specification moves discretion out of silent, continuous application into one published act of drafting. The translation table is that act.

Technical capacity doesn't determine whether a system gets scrutiny; institutions decide who may examine it, on what terms, and whether their judgment can slow deployment.

In 2020 the Royal Statistical Society offered Ofqual two independent statisticians for its grading model's technical advisory group. Ofqual's condition was a five-year NDA barring public comment; the nominees declined, and the model went ahead \autocite{rss2020ofqual}. It downgraded nearly 40\% of predicted grades, disproportionately at state schools with weaker histories, while raising grades disproportionately at smaller fee-paying schools \autocite{cnbc2020ofqual}. The results were reversed within days of the public outcry.

The expertise existed and was offered. The institution would accept it only on terms stripping reviewers of the ability to make their judgment public, and once the results produced consequences, public pressure achieved what pre-deployment review had been prevented from doing.
